% !TeX root = ../main.tex
\section{Type inference}
\label{sec:inference}

Writing the types in a program helps with design and documentation, but it
takes time. In a statically typed language the compiler can deduce the types
from the context in which each entity is used. In LETREC this means making the
annotations optional:
\begin{lstlisting}
Expression ::= letrec Type? Identifier (Identifier : Type?) = Expression in Expression
             | proc (Identifier : Type?) Expression
             | ...
Type       ::= int | boolean
\end{lstlisting}
The Rascal grammar of Section~\ref{sec:letrec} already allows a missing type;
the abstract syntax stores it as \texttt{noType()}.

\subsection{Rules with type variables}

When a type is unknown, the checker gives it a name, a \textbf{type variable}
$t_0, t_1, \dots$, and collects equations between types from the rules:
\begin{mathpar}
\inferrule*[right=(Diff)]{e_1 : \mathit{Int} \\ e_2 : \mathit{Int}}{\mathtt{-}(e_1, e_2) : \mathit{Int}}
\and
\inferrule*[right=(Pred)]{e_1 : \mathit{Int}}{\mathtt{zero?}(e_1) : \mathit{Boolean}}
\and
\inferrule*[right=(If)]{e_1 : \mathit{Boolean} \\ e_2 : T \\ e_3 : T}{(\mathtt{if}\ e_1\ \mathtt{then}\ e_2\ \mathtt{else}\ e_3) : T}
\and
\inferrule*[right=(Abs)]{v : T_{\mathit{var}} \\ \mathit{body} : T_{\mathit{body}}}{\mathtt{proc}(v)\ \mathit{body} : T_{\mathit{var}} \to T_{\mathit{body}}}
\and
\inferrule*[right=(Reduc)]{\mathit{rator} : T_1 \to T_2 \\ \mathit{rand} : T_1}{\mathtt{call}(\mathit{rator}, \mathit{rand}) : T_2}
\end{mathpar}
Every use of a rule produces equations. Solving the equations gives the
types.

\paragraph{Example 1.} The lecture starts with
\begin{lstlisting}
letrec foo (x) = if zero?(x) then 1 else -(x, (foo -(x,1))) in foo
\end{lstlisting}
Let $x : t_x$ and $\mathit{foo} : t_x \to t_r$. $\mathtt{zero?}(x)$ gives $t_x = \mathit{Int}$.
The \texttt{then} branch is $1 : \mathit{Int}$, so the \texttt{if} has type
$\mathit{Int}$ and $t_r = \mathit{Int}$. The \texttt{else} branch is a difference,
which needs $(\mathit{foo}\ \mathtt{-}(x,1)) : \mathit{Int}$, consistent with
$t_r = \mathit{Int}$. So $\mathit{foo} : \mathit{Int} \to \mathit{Int}$.

\paragraph{Example 2.} For
\begin{lstlisting}
proc (f) proc (x) -((f 3), (f x))
\end{lstlisting}
the lecture names every subexpression (Figure~\ref{fig:infer-tree}): $t_0$ for the
whole term, $t_1$ for the inner \texttt{proc}, $t_2$ for the difference, $t_3$
for \texttt{(f 3)}, $t_4$ for \texttt{(f x)}, and $t_f$, $t_x$ for the
parameters. The equations are in Table~\ref{tab:infer-eqs}.

\begin{figure}[h]
\centering
\begin{forest}
  for tree={font=\small, s sep=4mm, l sep=4mm, align=center},
  [{\texttt{proc (f)}\\$t_0$}
    [{\texttt{proc (x)}\\$t_1$}
      [{\texttt{-( , )}\\$t_2$}
        [{\texttt{(f 3)}\\$t_3$} [{\texttt{f}\\$t_f$}] [{\texttt{3}\\$\mathit{Int}$}]]
        [{\texttt{(f x)}\\$t_4$} [{\texttt{f}\\$t_f$}] [{\texttt{x}\\$t_x$}]]
      ]
    ]
  ]
\end{forest}
\caption{Type variables for the subexpressions of \texttt{proc (f) proc (x) -((f 3), (f x))}.}
\label{fig:infer-tree}
\end{figure}

\begin{table}[h]
\centering
\small
\begin{tabular}{l l l}
  \toprule
  Subexpression & Rule & Equation \\
  \midrule
  \texttt{(f 3)}         & Reduc & $t_f = \mathit{Int} \to t_3$ \\
  \texttt{(f x)}         & Reduc & $t_f = t_x \to t_4$ \\
  \texttt{-((f 3),(f x))} & Diff & $t_3 = \mathit{Int}$,\ \ $t_4 = \mathit{Int}$,\ \ $t_2 = \mathit{Int}$ \\
  \texttt{proc (x) \dots}  & Abs  & $t_1 = t_x \to t_2$ \\
  \texttt{proc (f) \dots}  & Abs  & $t_0 = t_f \to t_1$ \\
  \bottomrule
\end{tabular}
\caption{Equations for Example 2.}
\label{tab:infer-eqs}
\end{table}

Solving: $t_f = \mathit{Int} \to \mathit{Int}$; from the second equation
$t_x \to \mathit{Int} = \mathit{Int} \to \mathit{Int}$, so $t_x = \mathit{Int}$; then
$t_1 = \mathit{Int} \to \mathit{Int}$ and
\[
t_0 = (\mathit{Int} \to \mathit{Int}) \to \mathit{Int} \to \mathit{Int}.
\]
The function takes a function from integers to integers and returns a
function from integers to integers. The parentheses around the first
$\mathit{Int} \to \mathit{Int}$ matter: the slide writes
$t_0 = \mathit{Int} \to \mathit{Int} \to \mathit{Int} \to \mathit{Int}$, which, since the
arrow associates to the right, would be a function of three integers.

\paragraph{Example 3.} \texttt{if x then -(x,1) else 0}. The difference needs
$t_x = \mathit{Int}$; the condition of the \texttt{if} needs
$t_x = \mathit{Boolean}$. Together, $\mathit{Boolean} = \mathit{Int}$, which has no
solution: the expression has a type error.

\paragraph{Example 4.} \texttt{proc (f) zero?((f f))}. With $f : t_f$ and
$(f\ f) : t_3$, application gives $t_f = t_f \to t_3$, and \texttt{zero?} gives
$t_3 = \mathit{Int}$, so $t_f = t_f \to \mathit{Int}$. The type of $f$ would be defined
in terms of itself, and no finite type satisfies the equation. The
\textbf{occurs check} detects this: a variable cannot be bound to a type that
contains it.

\subsection{The algorithm}

The lecture divides type inference into five parts:
\begin{enumerate}
  \item substitution: replace a type variable by a known type inside a type
    expression;
  \item the no-occurrence property: a variable that is bound never appears
    in the types of the substitution, which the occurs check preserves;
  \item unification: add the equation $t_1 = t_2$ between two types to a
    substitution;
  \item fresh type variables for the expressions whose type is unknown;
  \item the type of an expression, computed from the rules and the current
    substitution.
\end{enumerate}

\paragraph{Substitution.} Substituting a type $t_1$ for a variable $t_{\mathit{var}}$
in a type expression $t_0$ replaces every occurrence:
\[
(t_{\mathit{var}} \to t_x \to t_y \to \mathit{Int} \to t_{\mathit{var}})[t_{\mathit{var}} = T]
\;=\; T \to t_x \to t_y \to \mathit{Int} \to T.
\]
The function goes by cases on $t_0$:
\begin{lstlisting}
substitution(t0, tvar, t1) : Type {
  switch (t0) {
    case int-type(): int-type;
    case boolean-type(): boolean-type;
    case proc-type(arg-type, result-type):
      proc-type(substitution(arg-type, tvar, t1), substitution(result-type, tvar, t1));
    case tvar-type(sn): if (t0 == tvar) then t1 else t0;
  }
}
\end{lstlisting}
(The slide compares \texttt{t1 = tvar} in the last case. The comparison is
between the variable found, $t_0$, and the variable being replaced.) A
substitution maps several variables to types. In Rascal, \texttt{apply} follows
chains such as $t_1 = t_2$, $t_2 = \mathit{Int}$:
\coderange{Rascal}{rascal/src/Infer.rsc}{subst}

\paragraph{Unification.} Given two types, the unifier returns a substitution in
which they are equal:
\begin{enumerate}
  \item apply the current substitution to $t_1$ and $t_2$;
  \item if the results are the same, return the substitution unchanged;
  \item if $t_1$ is a variable, propose $t_1 = t_2$, after the occurs check;
  \item if $t_2$ is a variable, do the same;
  \item if both are function types, unify the argument types and then the
    result types; otherwise ($\mathit{Int}$ against $\mathit{Boolean}$, or a basic
    type against a function type) there is no solution.
\end{enumerate}
\begin{lstlisting}
unifier(t1, t2, subs, exp) : Substitution {
  t1 = apply-substitution(t1, subs);
  t2 = apply-substitution(t2, subs);
  condition (t1 == t2) { subs; }
  tvar-type?(t1) {
    if (no-occurrence?(t1, t2)) extend-substitution(subs, t1, t2);
    else Error(No occurrence violation t1 t2 exp);
  }
  tvar-type?(t2) { ... }
  proc-type?(t1) && proc-type?(t2) {
    subs = unifier(proc-type.arg(t1), proc-type.arg(t2), subs, exp);
    subs = unifier(proc-type.res(t1), proc-type.res(t2), subs, exp);
  }
  else { Error(Unification failure t1 t2 exp); }
}

no-occurrence?(tvar, t) : Boolean {
  switch (t) {
    int-type: true;
    bool-type: true;
    proc-type(args, res): no-occurrence?(tvar, args) && no-occurrence?(tvar, res);
    var-type: not(tvar == t)
  }
}
\end{lstlisting}
In Rascal:
\coderange{Rascal}{rascal/src/Infer.rsc}{unify}

\paragraph{Types for expressions.} A missing annotation becomes a fresh
variable (\texttt{optionaltype-2-type}), and the inference function follows
the rules, calling the unifier at each premise that requires two types to be
equal. For an application it creates a variable for the result and unifies
the type of the operator with $t_{\mathit{rand}} \to t_{\mathit{res}}$:
\coderange{Rascal}{rascal/src/Infer.rsc}{infer}
The tests run the four examples of this section: \texttt{foo} gets
\texttt{int -> int}, Example 2 gets \texttt{(int -> int) -> int -> int}, Example 3
fails with a unification failure and Example 4 with a no-occurrence violation.

\subsection{Principal types and Algorithm W}

This subsection goes beyond the slides. The method above is a version of
Algorithm W, the inference algorithm of ML and Haskell (the Hindley--Milner
type system). Unification, when it succeeds, finds the \textbf{most general
unifier}: every other solution is an instance of it. So every typable term
has a \textbf{principal type}, of which all its other types are instances.
The principal type of \texttt{proc (x) x} is $t_0 \to t_0$; $\mathit{Int} \to \mathit{Int}$
is a valid type for it, but an instance. The Rascal test \texttt{identity}
shows the variable left in the result. Section~\ref{sec:polymorphism} uses this
to give \texttt{let}-bound functions polymorphic types.

Checking and inference are different problems. Checking verifies a type that
is given; inference has to find one. For richer type systems inference can
be undecidable, and many modern languages (Rust, Swift, TypeScript) use
\textbf{bidirectional} checking: some constructs synthesize their type
bottom-up (variables, literals, applications), and others are checked against
a type that comes from the context (an unannotated function passed where a
function type is expected). They need more annotations than ML, and combine
better with subtyping.

\paragraph{Generics in Java.} Java infers the type parameters at each call of a
generic method, a restricted form of the same problem:
\coderange{Java}{java/Generics.java}{generics}
\coderange{Java}{java/Generics.java}{use}
For \texttt{p} the compiler infers \texttt{MyPair<Integer, String>}, for
\texttt{q} and \texttt{result} \texttt{Integer}, and for \texttt{pairs}
\texttt{List<MyPair<Integer, String>>}.

\paragraph{References.} The same method extends to the commands of the
semantics block. A cell has type $\tref{T}$, and commands have type $\unit$:
\begin{mathpar}
\inferrule*[right=T-Deref]{\Gamma \vdash l : \tref{T}}{\Gamma \vdash \deref{l} : T}
\and
\inferrule*[right=T-Assign]{\Gamma \vdash l : \tref{T} \\ \Gamma \vdash e : T}{\Gamma \vdash \assign{l}{e} : \unit}
\and
\inferrule*[right=T-While]{\Gamma \vdash B : \tybool \\ \Gamma \vdash C : \unit}{\Gamma \vdash \whiledo{B}{C} : \unit}
\end{mathpar}
In the program
\[
\begin{array}{l}
\assign{a}{1}\ ;\ \assign{b}{1}\ ;\ \whiledo{\deref{n} < 5}{C} \\
C = \ift{\deref{\mathit{term}} \le 2}{\assign{\mathit{res}}{\deref{n}}}{(\assign{c}{\deref{a} + \deref{b}}\ ;\ \assign{\mathit{res}}{\deref{n}}\ ;\ \assign{a}{\deref{b}}\ ;\ \assign{b}{\deref{c}})}\ ;\ \assign{n}{\deref{n} + 1}
\end{array}
\]
every cell is dereferenced and compared with or added to an integer, so
$n, \mathit{res}, a, b, c, \mathit{term} : \tref{\tyint}$; the conditions are $\tybool$; and
every command, the \texttt{while} included, is $\unit$.

\subsection{Exercises}

\begin{exercise}
What is the inferred type of \texttt{if x then -(x,1) else 0}?
\end{exercise}

\paragraph{Solution.} None. The difference requires $t_x = \mathit{Int}$ and the
condition requires $t_x = \mathit{Boolean}$; unifying $\mathit{Int}$ with
$\mathit{Boolean}$ fails.

\begin{exercise}
Infer the type of \texttt{proc (f) proc (x) (f (f x))}.
\end{exercise}

\paragraph{Solution.} $(f\ x)$ gives $t_f = t_x \to t_1$ and $(f\ (f\ x))$ gives
$t_f = t_1 \to t_2$. Unifying the two function types: $t_x = t_1$ and $t_1 = t_2$.
So $t_f = t_x \to t_x$ and the type is $(t_x \to t_x) \to t_x \to t_x$: apply a function
twice, for any type.

\begin{exercise}
Infer the principal type of $\lambda x.\, \lambda y.\, \lambda z.\, x\, z\, (y\, z)$ (the S
combinator).
\end{exercise}

\paragraph{Solution.} Let $z : \gamma$. $y\, z$ gives $y : \gamma \to \beta$. $x\, z\, (y\, z)$
gives $x : \gamma \to \beta \to \alpha$. The type is
$(\gamma \to \beta \to \alpha) \to (\gamma \to \beta) \to \gamma \to \alpha$.

\begin{exercise}
Run the unifier by hand on $t_f = t_f \to \mathit{Int}$ and say at which step it stops.
\end{exercise}

\paragraph{Solution.} After applying the (empty) substitution the two types are
different, and $t_f$ is a variable, so the unifier proposes
$t_f = t_f \to \mathit{Int}$. The occurs check finds $t_f$ inside $t_f \to \mathit{Int}$ and
reports a no-occurrence violation. Without the check, the substitution would
replace $t_f$ forever.
